% Generado por src/comparativa.py — no editar a mano, se regenera.
\begin{table}[htbp]
    \centering
    {\footnotesize
    \begin{tabular}{@{}lrrrrrr@{}}
        \toprule
        \textbf{Observable} & \textbf{n (cobertura)} & \textbf{Ridge} & \textbf{RF} & \textbf{SAGE} & \textbf{GIN} & \textbf{GATv2} \\
        \midrule
        media de $Z$ & 1986 (50\%) & 0,3389 & 0,6370 & 0,5755 & 0,5317 & 0,4733 \\
        media de $X$ & 3393 (85\%) & 0,4978 & 0,7916 & 0,5500 & 0,7321 & 0,6180 \\
        media de $Y$ & 1096 (27\%) & 0,5164 & 0,6296 & 0,5693 & 0,5876 & 0,5283 \\
        paridad & 1825 (46\%) & 0,7282 & 0,8888 & 0,8373 & 0,9277 & 0,7858 \\
        corr. de vecinos & 3369 (84\%) & 0,3520 & 0,6079 & 0,5251 & 0,5830 & 0,4316 \\
        \bottomrule
    \end{tabular}}

    \caption{Proporción de circuitos en los que la corrección acerca el valor al exacto. Es la versión robusta de la mejora relativa. Conjunto de test, por observable. $n$ es el número de casillas evaluables. Ridge = Ridge de 34 variables, RF = Random Forest; SAGE, GIN y GATv2 son las tres convoluciones del MPNN.  Son métricas post-ejecución: se reportan, pero la comparativa la decide el $R^2$ sobre $f$.}
    \label{tab:mitigacion_fraccion}

    {\footnotesize Fuente: elaboración propia.}
\end{table}
